%% FullCourse_10Lecs -- Lecture 1, Topic 1.2: Equilibrium Concepts + Two Classical Approaches
%% Source: ./archive_pre_split/Lec01_Quant_Macro.tex (sections 3-4)
%% Pass A mechanical split (verbatim content preserved).

%% Shared preamble for the FullCourse_10Lecs topic decks.
%%
%% Each topic .tex \input's this file, then defines its own \title, \date
%% override (if any), \graphicspath, and \begin{document}.
%%
%% Reconstructed verbatim from the inlined copy preserved in
%% Lec10_Case_Studies/Lec10_T8_Case6_DeepRamsey.tex (lines 23-190), which is
%% explicitly annotated there as "from ../shared_preamble.tex, copied verbatim".
%% Base preamble matches MiniCourse_8hr/Slides/shared_preamble.tex; this version
%% additionally provides \sectiondivider, the conceptbox/cautionbox/examplebox/
%% takeawaybox tcolorboxes, and \compacttables used across the split decks.

\documentclass[10pt,english,aspectratio=169]{beamer}

\usepackage{ae,aecompl}
\usepackage[T1]{fontenc}
\usepackage{textcomp}
\usepackage[utf8]{inputenc}
\usepackage{babel}
\usepackage{datetime2}

\setcounter{secnumdepth}{3}
\setcounter{tocdepth}{3}

\usepackage{amsmath, amssymb, amsthm, graphicx}
\usepackage{booktabs, multirow}
\usepackage{tikz}
\usepackage{pgfplots}
\pgfplotsset{compat=1.16}
\usepackage[most]{tcolorbox}
\tcbuselibrary{skins, breakable, theorems}
\usepackage{listings}
\usepackage{xcolor}

\lstset{
  basicstyle=\ttfamily\small,
  keywordstyle=\color{blue!80!black}\bfseries,
  commentstyle=\color{green!50!black}\itshape,
  stringstyle=\color{orange!80!black},
  backgroundcolor=\color{gray!8},
  frame=single,
  rulecolor=\color{gray!40},
  breaklines=true,
  showstringspaces=false,
  numbers=none,
}

\lstdefinestyle{pythonstyle}{
    language=Python,
    basicstyle=\ttfamily\footnotesize,
    keywordstyle=\color{blue!80!black}\bfseries,
    commentstyle=\color{green!50!black}\itshape,
    stringstyle=\color{orange!80!black},
    frame=single,
    breaklines=true,
    showstringspaces=false,
    numbers=left,
    numbersep=5pt,
    numberstyle=\tiny\color{gray},
    backgroundcolor=\color{gray!8},
    rulecolor=\color{gray!40}
}

\ifx\hypersetup\undefined
  \AtBeginDocument{%
    \hypersetup{bookmarks=false,breaklinks=true,pdfborder={0 0 0},colorlinks=true,
      linkcolor=DarkText,urlcolor=RedTitle}%
  }
\else
  \hypersetup{bookmarks=false,breaklinks=true,pdfborder={0 0 0},colorlinks=true,
    linkcolor=DarkText,urlcolor=RedTitle}%
\fi

\makeatletter
\newcommand\makebeamertitle{%
  \begin{frame}[plain]
    \vspace*{1.0cm}
    \centering
    {\usebeamerfont{title}\usebeamercolor[fg]{title}\inserttitle\par}
    \vspace{1.0cm}
    {\usebeamerfont{author}\usebeamercolor[fg]{author}\insertauthor\par}
    \vspace{0.2cm}
    {\usebeamerfont{date}\usebeamercolor[fg]{date}\insertdate\par}
  \end{frame}}%
\AtBeginDocument{%
  \let\origtableofcontents=\tableofcontents
  \def\tableofcontents{\@ifnextchar[{\origtableofcontents}{\gobbletableofcontents}}
  \def\gobbletableofcontents#1{\origtableofcontents}%
}

\usetheme{metropolis}
\usefonttheme{professionalfonts}

\definecolor{LightGray}{RGB}{230,230,230}
\definecolor{RedTitle}{RGB}{0,0,200}
\definecolor{VeryLightGray}{RGB}{245,245,245}
\definecolor{DarkText}{RGB}{0,0,0}

\setbeamercolor{background canvas}{bg=white}
\setbeamercolor{title}{fg=RedTitle}
\setbeamercolor{author}{fg=DarkText}
\setbeamercolor{date}{fg=DarkText}
\setbeamercolor{frametitle}{bg=VeryLightGray,fg=DarkText}
\setbeamercolor{normal text}{fg=DarkText}
\setbeamerfont{title}{size=\large}

\setbeamersize{text margin left=5mm, text margin right=5mm}
\addtobeamertemplate{frametitle}{}{\vspace{-0.5em}}
\newcommand{\reducevspace}{\vspace{-0.5cm}}

% Shared section divider and box styles for split topic decks.
\newcommand{\sectiondivider}[2]{%
\begin{frame}[plain,noframenumbering]
\begin{center}
\vfill
{\Large\textbf{#1}\par}
\vspace{0.35cm}
{\normalsize #2\par}
\vfill
\end{center}
\end{frame}
}

\newtcolorbox{conceptbox}[1][]{
  colback=blue!3,
  colframe=RedTitle!55,
  boxrule=0.35mm,
  arc=1mm,
  left=2mm,
  right=2mm,
  top=1mm,
  bottom=1mm,
  #1
}
\newtcolorbox{cautionbox}[1][]{
  colback=red!3,
  colframe=red!50!black,
  boxrule=0.35mm,
  arc=1mm,
  left=2mm,
  right=2mm,
  top=1mm,
  bottom=1mm,
  #1
}
\newtcolorbox{examplebox}[1][]{
  colback=green!3,
  colframe=green!45!black,
  boxrule=0.35mm,
  arc=1mm,
  left=2mm,
  right=2mm,
  top=1mm,
  bottom=1mm,
  #1
}
\newtcolorbox{takeawaybox}[1][]{
  colback=VeryLightGray,
  colframe=RedTitle!55,
  boxrule=0.35mm,
  arc=1mm,
  left=2mm,
  right=2mm,
  top=1mm,
  bottom=1mm,
  #1
}
\newcommand{\compacttables}{\renewcommand{\arraystretch}{1.15}}

\setbeamertemplate{navigation symbols}{}
\setbeamertemplate{footline}{%
  \hfill%
  \usebeamercolor[fg]{page number in head/foot}%
  \usebeamerfont{page number in head/foot}%
  \insertframenumber\,/\,\inserttotalframenumber\kern1em\vskip2pt%
}

\makeatother

\author{Zhigang Feng}
% "date with time": compile timestamp so each build is identifiable.
\date{\today\ \DTMcurrenttime}


\graphicspath{{./pic/}{../../FullCourse_10Lecs/Lec01_Quant_Macro/pic/}{../}}

\title[AI for Econ Research]{AI for Economic Research: Dynamic Models, Language, and Agents\\
Lecture 1.2: Equilibrium Concepts \& Two Classical Approaches}

\begin{document}

\makebeamertitle

%=============================================================================
\begin{frame}{Topic 1.2 Agenda}
\begin{enumerate}
  \item \textbf{Equilibrium Concepts} --- competitive, recursive
  \medskip
  \item \textbf{Two Classical Approaches} --- VFI, Euler equation iteration
\end{enumerate}
\end{frame}

%=============================================================================
\section{Equilibrium Concepts}

\begin{frame}{Competitive Equilibrium: Definition}
\begin{itemize}
\item \textbf{A competitive equilibrium} consists of:
\begin{enumerate}
    \item \emph{Prices:} $\{r_t, w_t\}_{t=0}^\infty$ (rental rate and wage)
    \item \emph{Allocations:} $\{c_t, k_{t+1}, n_t\}_{t=0}^\infty$
\end{enumerate}
such that: (i) households optimize, (ii) firms optimize, (iii) markets clear.

\medskip
\item \textbf{Key Equilibrium Conditions:}
\begin{itemize}
    \item Factor prices: $r_t = \alpha A k_t^{\alpha-1} - \delta$, \quad $w_t = (1-\alpha)A k_t^\alpha$
    \item Household Euler equation: $u'(c_t) = \beta \mathbb{E}_t[u'(c_{t+1})(1 + r_{t+1})]$
    \item Resource constraint: $c_t + k_{t+1} = A k_t^\alpha + (1-\delta)k_t$
\end{itemize}

\medskip
\item \textbf{Recursive (Markov) Equilibrium:} An equilibrium where decisions depend only on \emph{payoff-relevant} state variables, not on calendar time or history.
\begin{itemize}
    \item Representative agent: state = $(k, z)$
    \item Policy rules: $c = \sigma(k,z)$, $k' = g(k,z)$, prices $r = r(k,z)$, $w = w(k,z)$
\end{itemize}
\end{itemize}
\end{frame}

\begin{frame}{From Competitive Equilibrium to Social Planner}
\begin{itemize}
\item \textbf{First Welfare Theorem:} Under complete markets and no externalities, a competitive equilibrium is Pareto optimal.

\medskip
\item \textbf{Implication for Computation:}
\begin{itemize}
    \item For representative agent models, solve the social planner's problem
    \item Recover prices from planner's marginal conditions
    \item Avoids fixed-point problem over prices
\end{itemize}

\medskip
\item \textbf{When the Theorem Fails:}
\begin{itemize}
    \item Incomplete markets (Aiyagari, Bewley): no Pareto optimality
    \item Externalities, taxes, monopolistic competition
    \item Must solve for equilibrium prices \emph{explicitly}
\end{itemize}

\medskip
\item \textbf{Practical Implication:}
\begin{itemize}
    \item Representative agent models: ``easy'' (solve planner's problem)
    \item Heterogeneous agent models: ``hard'' (must iterate on distribution + prices)
\end{itemize}
\end{itemize}
\end{frame}

\begin{frame}{Two Roads to Equilibrium}
\begin{itemize}
\item We have two constructive algorithms for computing equilibrium:

\bigskip
\item \textbf{Approach 1: Dynamic Programming (Value Function Iteration)}
\begin{itemize}
    \item Iterate on the Bellman equation: $v_{n+1} = Tv_n$
    \item Convergence via \emph{contraction mapping}
    \item Recovers both value function and policy
\end{itemize}

\bigskip
\item \textbf{Approach 2: Euler Equation Iteration (Greenwood-Huffman)}
\begin{itemize}
    \item Iterate directly on the policy function via Euler equation
    \item Convergence via \emph{monotonicity}
    \item Works directly with equilibrium conditions
\end{itemize}

\bigskip
\item Both approaches are \emph{constructive}---they provide algorithms, not just existence proofs.
\end{itemize}
\end{frame}

%=============================================================================
\section{Two Classical Approaches}

\begin{frame}{Dynamic Programming: The Classical Framework}
\begin{itemize}
\item \textbf{Dynamic Programming (Bellman, 1957):} Solve sequential decision problems by exploiting recursive structure.

\medskip
\item \textbf{Core Idea:} Value of any state = max over (immediate reward + discounted continuation):
\[
v(s) = \max_{a \in \Gamma(s)} \left\{ r(s, a) + \beta \, \mathbb{E}\left[ v(s') \mid s, a \right] \right\}
\]

\medskip
\item \textbf{Value Function Iteration (VFI):}
\begin{enumerate}
    \item Discretize state space into finite grid $\{s_1, \ldots, s_N\}$
    \item Initialize $v_0(s_i)$ for each grid point
    \item Iterate: $v_{n+1}(s_i) = \max_{a} \{ r(s_i, a) + \beta \sum_j P(s_j|s_i,a) v_n(s_j) \}$
    \item Stop when $\|v_{n+1} - v_n\| < \varepsilon$
\end{enumerate}

\medskip
\item \textbf{Requirement:} Complete knowledge of the environment---$P(s'|s,a)$ and $r(s,a)$.
\end{itemize}
\end{frame}

\begin{frame}{Euler Equation Approach: Setup}
\begin{itemize}
\item \textbf{Alternative:} Work directly with first-order conditions (Euler equations).

\medskip
\item The Euler equation for the optimal growth model:
\[
u'(c_t) = \beta \, \mathbb{E}_t \left[ u'(c_{t+1}) \cdot F_k(k_{t+1}, z_{t+1}) \right]
\]

\medskip
\item In a competitive equilibrium with $k = K$:
\begin{multline*}
u_c(c(k,K,k',z)) = \\
\beta \, \mathbb{E} \left[ u_c(c(k',K',k'',z')) \cdot F_k(k',K',z') \mid z \right]
\end{multline*}

\medskip
\item \textbf{Equilibrium Requirement:} Find policy $g$ and aggregate law of motion $G$ such that:
\begin{itemize}
    \item The Euler equation holds for all $(k, K, z)$
    \item Consistency: $G(K, z) = g(K, K, z)$ (individual = aggregate)
\end{itemize}
\end{itemize}
\end{frame}

\begin{frame}{The Greenwood-Huffman Operator}
\begin{itemize}
\item \textbf{Idea:} Construct a sequence of policy functions that converges to equilibrium.

\medskip
\item Let $H^{0}(K,z) = 0$. Define $H^{j+1}(K,z)$ as the $x$ solving:
\[
u_c(c(K,K,x,z)) = \beta \, \mathbb{E} \left[ u_c(c(x,x,H^{j}(x,z'),z')) \cdot F_k(x,x,z') \right]
\]

\medskip
\item \textbf{Interpretation:}
\begin{itemize}
    \item Given a conjecture $H^j$ for future policy, solve for today's optimal $x$
    \item The solution $x = H^{j+1}(K,z)$ becomes the new conjecture
    \item Iterate until convergence: $H^{j+1} \approx H^j$
\end{itemize}

\medskip
\item \textbf{Convergence via Monotonicity:}
\begin{itemize}
    \item $\{H^j\}$ is monotone (increasing from $H^0 = 0$)
    \item Bounded by resource constraint $H^j(K,z) \leq F(K,K,z)$
    \item Monotone + bounded $\Rightarrow$ limit $H^*$ exists
\end{itemize}
\end{itemize}
\end{frame}

\begin{frame}{Comparing the Two Approaches}
\begin{center}
\renewcommand{\arraystretch}{1.3}
\begin{tabular}{|l|c|c|}
\hline
 & \textbf{Value Function Iteration} & \textbf{Euler Equation Iteration} \\
\hline
\textbf{Object} & Value function $v(k)$ & Policy function $H(K,z)$ \\
\textbf{Iteration} & $v_{n+1} = Tv_n$ & $H^{j+1} = TH^j$ \\
\textbf{Convergence} & Contraction & Monotonicity \\
\textbf{Theory} & Banach fixed point & Monotone convergence \\
\hline
\textbf{Strength} & General, robust & Direct equilibrium \\
\textbf{Weakness} & Indirect (need to extract $g$) & Requires monotonicity \\
\hline
\end{tabular}
\end{center}

\bigskip
\textbf{Both provide:}
\begin{itemize}
    \item Constructive algorithms (not just existence proofs)
    \item Theoretical guarantees for convergence
    \item Foundation for numerical implementation
\end{itemize}
\end{frame}

\begin{frame}{Why Classical Macro Models Are ``Easy''}
\begin{itemize}
\item \textbf{Representative Agent + First Welfare Theorem:}
\begin{itemize}
    \item Competitive equilibrium $\Leftrightarrow$ Social planner's solution
    \item No need to explicitly compute prices to find allocations
\end{itemize}

\medskip
\item \textbf{The Key Simplification:} Individual capital $k$ = Aggregate capital $K$
\begin{itemize}
    \item Prices are deterministic functions of a single state variable:
    \[
    r(K) = \alpha A K^{\alpha - 1} - \delta, \qquad w(K) = (1-\alpha) A K^{\alpha}
    \]
    \item The agent \emph{knows} how prices evolve because $K' = g(K)$
\end{itemize}

\medskip
\item \textbf{By Design, Not by Accident:}
\begin{itemize}
    \item Economists constructed this framework for tractability
    \item The ``environment'' is fully specified by $(A, \alpha, \beta, \delta)$
    \item State space is low-dimensional (often just $k$ or $(k, z)$)
\end{itemize}
\end{itemize}
\end{frame}

\end{document}
